% appendices/appendix_math.tex
% Version 2.0 — Expanded to support self-contained reading
\chapter{最低必要数学背景}
\label{app:math}

本附录汇集了理解本书所需的核心数学工具。
目标不是完整的线性代数或优化教材，
而是让读者在遇到具体推导时可以快速查阅。
最重要的提示：\textbf{如果你理解了外积和矩阵梯度，
本书 80\% 的公式就会变得自然}。

%═════════════════════════════════════════════
\section{线性代数核心}
%═════════════════════════════════════════════

\subsection{外积（Outer Product）}

向量 $a \in \R^m, b \in \R^n$ 的外积为矩阵：
\[
  a b^\top \in \R^{m \times n}, \quad (ab^\top)_{ij} = a_i b_j
\]

\textbf{核心性质}：
\begin{itemize}
  \item $ab^\top$ 是一个\textbf{秩 1 矩阵}——它的所有列都是 $a$ 的标量倍。
  \item 对任意向量 $c$：$(ab^\top) c = a (b^\top c)$——
    右乘一个向量，结果是 $a$ 方向上的投影，幅度由 $b^\top c$ 决定。
\end{itemize}

\textbf{为什么外积在本书中反复出现}：
\begin{enumerate}
  \item \textbf{联想记忆}（第~\ref{ch:fastweights}~章）：
    存储一对 $(k, v)$ 到矩阵中：$\mS \leftarrow \mS + v k^\top$。
    读取时 $\mS k' = v(k^\top k')$——如果 $k' \approx k$，输出 $\approx v$。
  \item \textbf{梯度结构}：损失 $\ell = \|Wx - y\|^2$ 的梯度 $\nabla_W \ell$
    也是外积形式（见下一节）。
  \item 这就是为什么``快权重更新''和``梯度步''在数学上长得一样。
\end{enumerate}

\subsection{矩阵的 Frobenius 范数}

$\|A\|_F = \sqrt{\sum_{i,j} A_{ij}^2} = \sqrt{\text{tr}(A^\top A)}$。

在本书中，正则化项 $\frac{\gamma}{2}\|\mS\|_F^2$
惩罚状态矩阵的``总能量''，防止无限增长。

%═════════════════════════════════════════════
\section{最小二乘与岭回归}
\label{sec:app_regression}
%═════════════════════════════════════════════

\subsection{普通最小二乘（OLS）}

给定数据矩阵 $X \in \R^{n \times d}$（$n$ 个样本，$d$ 维特征）
和目标 $Y \in \R^{n \times m}$，OLS 求解：
\[
  W^* = \arg\min_W \|XW - Y\|_F^2
\]

\textbf{正规方程（Normal Equations）}：
\begin{equation}
  W^* = (X^\top X)^{-1} X^\top Y
  \label{eq:ols_solution}
\end{equation}

\textbf{几何解释}：$XW^*$ 是 $Y$ 在 $X$ 列空间上的正交投影。
$(X^\top X)^{-1} X^\top$ 是伪逆矩阵 $X^+$。

当 $X^\top X$ 不可逆（$n < d$ 或特征共线性）时，OLS 不唯一。

\subsection{岭回归（Ridge Regression）}

加入 $L_2$ 正则化解决 $X^\top X$ 的不可逆问题：
\begin{equation}
  W^* = \arg\min_W \|XW - Y\|_F^2 + \gamma \|W\|_F^2
  = (X^\top X + \gamma I)^{-1} X^\top Y
  \label{eq:ridge_solution}
\end{equation}

$\gamma > 0$ 保证 $X^\top X + \gamma I$ 总是可逆的。
在本书中，第~\ref{ch:unified}~章的 TTR 目标函数
正是在线版本的岭回归。

%═════════════════════════════════════════════
\section{递归最小二乘（RLS）}
\label{sec:app_rls}
%═════════════════════════════════════════════

在在线设定中，数据点逐个到来。
目标：每次新增一个样本 $(x_t, y_t)$ 时，
不重新计算整个 $(X^\top X)^{-1}$，
而是用 $O(d^2)$ 的代价增量更新。

\subsection{推导思路}

设到时刻 $t-1$ 时，我们维护：
\begin{itemize}
  \item 状态矩阵 $W_{t-1} = P_{t-1} \cdot (\text{累积} X^\top Y)$
  \item 精度矩阵 $P_{t-1} = (X_{1:t-1}^\top X_{1:t-1} + \gamma I)^{-1}$
\end{itemize}

新增样本 $(x_t, y_t)$ 后：
\[
  P_t^{-1} = P_{t-1}^{-1} + x_t x_t^\top
\]

这是一个\textbf{秩 1 更新}。
用 Sherman-Morrison 公式求 $P_t$（见下一节），
避免了 $O(d^3)$ 的矩阵求逆。

\subsection{RLS 更新公式}

\begin{equation}
  \kappa_t = \frac{P_{t-1} x_t}{1 + x_t^\top P_{t-1} x_t}
  \label{eq:rls_gain}
\end{equation}
\begin{equation}
  W_t = W_{t-1} + (y_t - W_{t-1} x_t) \kappa_t^\top
  \label{eq:rls_update}
\end{equation}
\begin{equation}
  P_t = P_{t-1} - \kappa_t x_t^\top P_{t-1}
  \label{eq:rls_precision}
\end{equation}

这里 $\kappa_t$ 称为\textbf{Kalman 增益}。
$y_t - W_{t-1} x_t$ 是预测残差（prediction error）。

\textbf{观察}：
式~\eqref{eq:rls_update} 的形式是
$W_t = W_{t-1} + \text{(残差)} \cdot \text{(方向)}^\top$——
又是一个外积更新！这与 TTT-Layer 的梯度步在结构上完全一致。

%═════════════════════════════════════════════
\section{Sherman-Morrison 公式}
\label{sec:app_sm}
%═════════════════════════════════════════════

\begin{theorem}[Sherman-Morrison]
设可逆矩阵 $A \in \R^{d \times d}$，向量 $u, v \in \R^d$，
且 $1 + v^\top A^{-1} u \neq 0$，则：
\begin{equation}
  (A + uv^\top)^{-1} = A^{-1} - \frac{A^{-1} u v^\top A^{-1}}{1 + v^\top A^{-1} u}
  \label{eq:sherman_morrison}
\end{equation}
\end{theorem}

\subsection*{推导要点}

设 $B = A + uv^\top$，令 $B^{-1} = A^{-1} + C$（待定修正项）。

由 $B B^{-1} = I$：
\begin{align}
  (A + uv^\top)(A^{-1} + C) &= I \\
  I + AC + uv^\top A^{-1} + uv^\top C &= I \\
  AC &= -u v^\top A^{-1} - u v^\top C \\
  AC &= -u(v^\top A^{-1} + v^\top C)
\end{align}

猜测 $C$ 的形式为 $C = \alpha \cdot A^{-1} u v^\top A^{-1}$（秩 1），
代入求解 $\alpha = -\frac{1}{1 + v^\top A^{-1} u}$，即得证。

\textbf{核心用途}：
在 RLS 中，$A = P_{t-1}^{-1}$，$u = v = x_t$。
每次新增一个样本，只需 $O(d^2)$ 的矩阵向量乘法
就能更新逆矩阵，无需 $O(d^3)$ 求逆。

%═════════════════════════════════════════════
\section{矩阵微积分基础}
\label{sec:app_matrix_calc}
%═════════════════════════════════════════════

\subsection{标量对矩阵的梯度}

设 $L(W) = \frac{1}{2}\|Wx - y\|_2^2$，其中 $W \in \R^{m \times d}$。

展开：
\begin{align}
  L(W) &= \frac{1}{2}(Wx - y)^\top(Wx - y) \\
  \nabla_W L &= (Wx - y) x^\top \in \R^{m \times d}
  \label{eq:matrix_gradient}
\end{align}

这是一个\textbf{外积}：误差向量 $(Wx - y)$ 乘以输入 $x^\top$。
这就是为什么本书中梯度步和快权重更新在代数上完全一样。

\subsection{链式法则（矩阵版本）}

设 $L = g(f(W))$，其中 $f: \R^{m \times d} \to \R^{p}$，$g: \R^p \to \R$。
则：
\[
  \frac{\partial L}{\partial W_{ij}} = \sum_k \frac{\partial g}{\partial f_k} \cdot \frac{\partial f_k}{\partial W_{ij}}
\]

在 MAML（第~\ref{ch:bilevel}~章）中，
外循环梯度需要对``内循环结果相对于初始化参数''求导，
这涉及的链式法则包含 Hessian 矩阵（见下一小节）。

\subsection{Jacobian 与 Hessian}

\textbf{Jacobian}：$J = \frac{\partial f}{\partial x} \in \R^{m \times d}$，
函数 $f: \R^d \to \R^m$ 的一阶导数矩阵。
第 $i$ 行 $j$ 列是 $\frac{\partial f_i}{\partial x_j}$。

\textbf{Hessian}：$H = \frac{\partial^2 L}{\partial x^2} \in \R^{d \times d}$，
标量函数 $L: \R^d \to \R$ 的二阶导数矩阵。

\subsection{预处理（Preconditioning）}

标准梯度下降 $x \leftarrow x - \eta \nabla L$ 在高曲率方向上振荡。
预处理梯度下降用 $H^{-1}$（或其近似）来校正：
\[
  x \leftarrow x - \eta H^{-1} \nabla L
\]

这使得高曲率方向上的步长自动缩小，低曲率方向上自动放大。
在第~\ref{ch:icl}~章中，Transformer 的 FFN 被解释为
在近似计算 $H^{-1}$，从而实现了隐式的预处理梯度下降。

%═════════════════════════════════════════════
\section{在线学习与遗憾（Regret）}
\label{sec:app_online}
%═════════════════════════════════════════════

\subsection{设定}

在在线学习中，学习器与数据流交互：
\begin{enumerate}
  \item 时步 $t$：学习器做预测 $\hat{y}_t$
  \item 环境揭示真实值 $y_t$
  \item 学习器承受损失 $\ell_t(\hat{y}_t)$ 并更新策略
\end{enumerate}

\subsection{累积遗憾}

\begin{equation}
  \text{Regret}(T) = \sum_{t=1}^T \ell_t(\hat{y}_t)
  - \min_{y^*} \sum_{t=1}^T \ell_t(y^*)
  \label{eq:regret}
\end{equation}

第一项是学习器的累积损失，第二项是``事后诸葛亮最优''的累积损失。
好的在线算法的遗憾增长率为 $O(\sqrt{T})$（次线性），
意味着平均每步遗憾趋向零。

\subsection{Worked Example：Follow-the-Leader}

设 $\ell_t(w) = (w - y_t)^2$，标量情形。
Follow-the-Leader (FTL) 策略：$\hat{w}_t = \frac{1}{t-1}\sum_{i=1}^{t-1} y_i$（历史均值）。

对于 $T$ 步，FTL 的遗憾为 $O(\log T)$——
非常好，因为 $\log T \ll T$。

\textbf{与本书的联系}：
线性注意力的外积累积 $\mS_t = \sum_{i \leq t} v_i k_i^\top$
可以被视为一种加权的 Follow-the-Leader 策略。
TTT-Layer 的在线 SGD 更新
对应于 Online Gradient Descent，其遗憾为 $O(\sqrt{T})$。

%═════════════════════════════════════════════
\section{常用符号对照表}
\label{sec:app_notation}
%═════════════════════════════════════════════

\begin{center}
\begin{tabular}{lll}
\toprule
\textbf{符号} & \textbf{含义} & \textbf{首次出现} \\
\midrule
$\theta$ & 慢权重 / 全局参数 & 第~\ref{ch:intro}~章 \\
$\mS_t$ & 快权重 / 隐状态矩阵 & 第~\ref{ch:fastweights}~章 \\
$W_t^{\text{inner}}$ & 内层学习器参数 & 第~\ref{ch:ttt_bridge}~章 \\
$k_t, v_t, q_t$ & 键、值、查询向量 & 第~\ref{ch:attention}~章 \\
$\eta$ & 学习率 & 第~\ref{ch:bilevel}~章 \\
$\ell^{(l)}$ & 第 $l$ 层局部目标函数 & 第~\ref{ch:nested}~章 \\
$T^{(l)}$ & 第 $l$ 层的时间尺度 & 第~\ref{ch:nested}~章 \\
$\lambda$ & 遗忘因子 / 衰减率 & 第~\ref{ch:attention}~章 \\
\bottomrule
\end{tabular}
\end{center}
